\documentclass[11pt]{article}
\usepackage{amsmath,amssymb,amsthm}
\usepackage[hidelinks]{hyperref}

% Reproducible pdfLaTeX metadata. The paper is compiled with shell escape off.
\ifdefined\pdfsuppressptexinfo\pdfsuppressptexinfo=-1\fi
\ifdefined\pdftrailerid\pdftrailerid{}\fi
\hypersetup{
  pdftitle={An Exact Counterexample to Graffiti 197},
  pdfauthor={AdaIvy project},
  pdfsubject={Exact Kneser K(7,2) counterexample with linked Lean certificate}
}

\setlength{\textwidth}{6.35in}
\setlength{\textheight}{8.9in}
\setlength{\oddsidemargin}{0.05in}
\setlength{\evensidemargin}{0.05in}
\setlength{\topmargin}{-0.35in}
\setlength{\parindent}{0pt}
\setlength{\parskip}{0.65em}
\setlength{\emergencystretch}{2em}

\newtheorem{proposition}{Proposition}

\title{An Exact Counterexample to Graffiti 197}
\author{AdaIvy project}
\date{21 August 2026}

\begin{document}
\maketitle

\begin{center}
\fbox{\parbox{0.88\textwidth}{\textbf{Status.}
This note gives a human-checkable exact counterexample. This result was
AI-generated by the AdaIvy project. Its closed arithmetic certificate was
kernel-checked in Lean with no axioms; the spectral derivation itself remains
human-checkable rather than fully formalized. It creates no AdaIvy epistemic
warrant, and novelty, significance, and publication approval remain
unassessed.}}
\end{center}

\section{Statement and conventions}

The conjecture known as Graffiti 197 states
\cite[Graffiti 197, PDF p.~71]{fajtlowicz2004}
\[
  -\lambda_{\mathrm{second\;smallest}}(A(G))
  \;\leq\;
  \operatorname{range}(\operatorname{spec}(\Gamma(G))),
\]
where $A(G)$ is the adjacency matrix and $\Gamma(G)$ is the gravity matrix.
This note uses the source-supported convention that the range of a vector is
the number of its distinct values.

For a graph with $n$ vertices, degrees $d(u)$, and graph distance
$\operatorname{dist}(u,v)$, the gravity matrix is
\[
 \Gamma_{uv}=
 \begin{cases}
 \displaystyle\frac{d(u)d(v)}{(n-1)\operatorname{dist}(u,v)},&u\ne v,\\[6pt]
 0,&u=v.
 \end{cases}
\]

\begin{proposition}
Graffiti 197 is false under these conventions.  The Kneser graph
$K(7,2)$ is a counterexample: its left-hand side is $4$, while its gravity
matrix has exactly $3$ distinct eigenvalues.
\end{proposition}

\paragraph{Auditable derivation.}
The investigation used the Kneser-graph candidate reported by AI Village News
\cite{aivillage2026} as a lead, then independently reconstructed the graph over
exact integers and rationals. The argument below counts the strongly regular
parameters, derives the adjacency spectrum, transforms it into the gravity
spectrum, and compares the two sides. It is a reproducible mathematical account,
not a claim to expose private model chain-of-thought.

\section{The graph and its adjacency spectrum}

The vertices of $K(7,2)$ are the two-element subsets of a seven-element set;
two vertices are adjacent exactly when the corresponding subsets are disjoint.
Consequently
\[
 n=\binom72=21,
 \qquad
 k=\binom52=10.
\]

The graph is strongly regular with parameters
\[
 (n,k,\lambda,\mu)=(21,10,3,6).
\]
Indeed, two adjacent two-subsets use four points and have
$\binom32=3$ common neighbors.  Two distinct nonadjacent two-subsets share one
point, use three points, and have $\binom42=6$ common neighbors.  Hence
\[
 A^2=(k-\mu)I+(\lambda-\mu)A+\mu J
     =4I-3A+6J.
\]

The all-ones vector has adjacency eigenvalue $10$.  On its orthogonal
complement $J$ vanishes, so every remaining adjacency eigenvalue $x$ obeys
\[
 x^2+3x-4=0,
 \qquad\text{hence}\qquad x\in\{1,-4\}.
\]
Writing their multiplicities as $f$ and $g$, dimension and trace give
\[
 f+g=20,
 \qquad
 10+f-4g=0,
\]
and therefore $f=14$ and $g=6$.  Thus
\[
 \operatorname{spec}(A)=\{10^1,1^{14},(-4)^6\}.
\]
Because $-4$ has multiplicity six, the second-smallest adjacency eigenvalue is
$-4$.  The left-hand side of Graffiti 197 is therefore
\[
 -\lambda_{\mathrm{second\;smallest}}(A)=4.
\]

\section{The gravity spectrum}

The graph has diameter two.  Since every degree is $10$ and $n-1=20$,
distinct adjacent vertices have gravity entry $100/20=5$, while distinct
nonadjacent vertices have entry $100/40=5/2$.  Therefore
\[
 \Gamma=\frac52(A+J-I).
\]

On the all-ones vector,
\[
 \Gamma\mathbf{1}
 =\frac52(10+21-1)\mathbf{1}
 =75\mathbf{1}.
\]
On the orthogonal complement of $\mathbf{1}$, the matrix $J$ vanishes.  The
adjacency eigenvalue $1$ is sent to
\[
 \frac52(1-1)=0,
\]
and the adjacency eigenvalue $-4$ is sent to
\[
 \frac52(-4-1)=-\frac{25}{2}.
\]
It follows that
\[
 \operatorname{spec}(\Gamma)
 =\left\{75^1,0^{14},\left(-\frac{25}{2}\right)^6\right\}.
\]
The gravity spectrum has exactly three distinct values.  Graffiti 197 would
therefore require
\[
 4\leq3,
\]
which is false.  This completes the counterexample calculation.

\section{Verification and scope}

An independent Python standard-library checker reconstructed $K(7,2)$, formed
the adjacency and gravity matrices over the rational numbers, computed their
characteristic polynomials exactly, and checked the stated factorizations.
The checker uses neither floating point nor the downloaded candidate program.

\paragraph{Run disclosure.}\mbox{}\par\noindent
AdaIvy's durable ledger records three completed OpenAI calls to
\nolinkurl{gpt-5-mini-2025-08-07}. The recorded external model cost is USD
0.017859 against the authorized USD 20 cap. Recorded usage is 7,747 input
tokens and 7,960 output tokens, for 15,707 tokens total. This
accounting is partial: the surrounding interactive Codex authoring session did
not expose billable usage to AdaIvy's ledger, so its tokens and cost are not
included or estimated.

AdaIvy also generated a closed Lean arithmetic certificate covering the four
combinatorial counts, the three gravity-eigenvalue transformations, and the
final contradiction $\neg(4\leq3)$. The sealed Lean 4.32.1 checker accepted it
with no axioms. The exact checked source is linked here:
\href{run:graffiti-197-kneser-counterexample.lean}%
{\nolinkurl{graffiti-197-kneser-counterexample.lean}}.

\begin{itemize}
\item Exact checker SHA-256:
\par{\scriptsize\texttt{2d487616404964fa46306fefc4ec7a0156e7d267377ee8d0a35280bd0ce4ad39}}.
\item Exact result content hash:
\par{\scriptsize\texttt{sha256:626439796ee4635919cdc3dc8c08f6f7bc13f8928770c7d15b8134905b090a1f}}.
\item Lean finding:
\texttt{formal-finding.db32ea12d524eaf9078e94ae}; outcome
\texttt{kernel\_checked}; approved and unapproved axiom lists both empty.
\item Lean wrapper SHA-256:
\par{\scriptsize\texttt{sha256:24fb637e1b7f06b656392636b85fdc1d5097ee5430c986021721bce0c7fce249}}.
\item Human-review status: this mathematical result has not yet been approved by
a human reviewer in AdaIvy. Automated checks alone do not make it an endorsed
result (machine status: \texttt{awaiting\_human\_review}; no warrant created).
\end{itemize}

The Lean certificate does not prove from first principles that the explicit
$21\times21$ adjacency and gravity matrices have the displayed spectra. It
checks the closed exact arithmetic consequences used after that mathematical
derivation. A complete matrix-spectrum formalization would require a broader
formal representation and approved imports. This result does not settle
Graffiti 322. If ``range'' is instead interpreted as numerical maximum minus
minimum, that is a different formulation; Roucairol and Cazenave report a
numerical $C_{17}$ counterexample for it \cite{roucairol2026}.

\begin{thebibliography}{9}
\bibitem{fajtlowicz2004}
S. Fajtlowicz, \emph{Written on the Wall}, July 2004 compilation,
Graffiti 197 on PDF page 71.
\href{https://independencenumber.wordpress.com/wp-content/uploads/2012/08/wow-july2004.pdf}%
{Source PDF}.

\bibitem{roucairol2026}
M. Roucairol and T. Cazenave, \emph{Refutation of Spectral Graph Theory
Conjectures with Search Algorithms}, arXiv:2409.18626v1 (2024); published in
CCIS 2976 (2026). \url{https://doi.org/10.1007/978-3-032-28263-7_2}.

\bibitem{aivillage2026}
AI Village News, \emph{Opus 5 \#124 --- WOW 197 gravity range FALSE},
13 August 2026. Treated as an unreviewed candidate source and independently
rechecked here.
\end{thebibliography}

\par\bigskip\noindent\emph{About this project}\footnote{AdaIvy is an open-source
AI-assisted mathematical research project that keeps retrieval, computation,
formal checking, human review, and publication status as separate recorded
boundaries. Source code and reproducibility materials:
\url{https://github.com/jketts/AdaIvy}.}

\end{document}
